\emph{The slice at \(k\ne0\): refinement of the window and model-level growth control.} On the slice \(k^0\mathrm e^{-3\pi\mathrm i/8}=r^0\in\mathbb R\), \(k^i\mathrm e^{\pi\mathrm i/8}=r^i\in\mathbb R\) the Cauchy-Binet identity \(\rho_X^2=\sum_{a=1}^{2}{(k^\mu A^a_\mu)}^2\) is exact at \(N=2\), and on \(\mathcal R_1\) the contraction \(k^\mu A_\mu\) reduces to the Hermitian real-linear form \(r^\mu\tilde A_\mu\), so \(\rho_X\) is real and \(|O_k|\le2\) pointwise there; right-hand-side boundedness on the slice is structural. The commuting-ray alternative also refines inside the slice: at \(u=0\), writing the growth functional as \(B=\mathbf v\mathbf v^{\mathrm T}-\mathbf u\mathbf u^{\mathrm T}\) with \(\mathbf u=(\cos(3\pi/8)r^0,\cos(\pi/8)\vec r\,)\), \(\mathbf v=(\sin(3\pi/8)r^0,-\sin(\pi/8)\vec r\,)\), the matrix \(B\) is supported on \(\operatorname{span}\{\mathbf e_0,\hat{\mathbf r}\}\), has rank at most two, and \(\lambda_{\max}(B)=[{(r^0)}^2-r^2+\sqrt{{(r^0)}^4+6{(r^0)}^2r^2+r^4}]/(2\sqrt2)\), which is positive iff \(r^0\ne0\); for pure spatial slice momenta (\(r^0=0\)) there is no commuting ray with growth, the neutral set is \(\hat{\mathbf r}\cdot\hat\rho=0\), and \(B\) is negative off it. For \(r^0\ne0\) the one-dimensional branch integral \(\int_{\mathbb R}\mathrm e^{-Qa^2}\,2\cos(\gamma_{\mathrm c}|a|)\,\mathrm da=2\sqrt{\pi/Q}\,\mathrm e^{-\gamma_{\mathrm c}^2/(4Q)}\), \(\gamma_{\mathrm c}=\mathrm e^{3\pi\mathrm i/8}r^0/2\), gives pointwise modulus growth \(\mathrm e^{{(r^0)}^2\cos(3\pi u/4-\pi/4)/(16\epsilon)}\) for the soft temporal mode at \(\gamma=0\), saturating as \(\epsilon\to0^+\) at fixed \(\gamma<0\) to a constant of order \(\mathrm e^{{(r^0)}^2\sin(3\pi u/4-\pi/4)/(8|\gamma|)}\) that diverges as \(\gamma\to0^-\). In the reduced \(N=2\) integral (regulator (A), contours \(\mathcal R_u\), \(\epsilon\to0^+\) last) with this insertion at pure temporal slice momenta, the modulus ratio \(|Z_\epsilon[O_k]|/|Z_\epsilon[1]|\) stays of order one and tends to \(2\) along the ladder down to \(\epsilon=3\cdot10^{-3}\) (fit slope \(0.0000\) of \(\ln\)-ratio against \(1/\epsilon\), against the pointwise corner value \(0.044/\epsilon\) at \(r^0=1\)), and at \(u=0.6\) with \(\gamma=-0.2,-0.05\) it stays between \(0.75\) and \(2.2\): in the reduction neither the commuting-ray growth \(\mathrm e^{c_2(k)/\epsilon}\) nor the pointwise saturation factor materialises, the corner branch growth being phase-space suppressed near the soft corner by Vandermonde and measure powers. These are reduced-model computations with oscillation-limited resolution at \(\gamma=0\); whether the left-hand side is controlled on the slice at \(k\ne0\) with \(r^0\ne0\) remains open at full \(N=2\), the \(\gamma\to0^-\) bridge from fixed-\(\gamma<0\) success remains open, and the pure-spatial half of the slice is now ray-free; Problem~E remains open and PROGRAMME-CONDITIONAL.
